% dynamics/runtime_dynamics.tex

\chapter{Runtime Dynamics}

\label{chap:runtime-dynamics}

% ============================================================
\section*{Motivation}
% ============================================================

Railway alignments are not merely stored geometric objects.
Within AIM, they are interpreted as state-generating structures.

The preceding chapters introduced railway states and vehicle-relative
interpretation.  This chapter now introduces the operational layer that
describes how runtime states evolve.

Geometry defines what the railway is.

Runtime dynamics define which operational states emerge from it.

\begin{principle}[Runtime-dynamics principle]
\label{princ:runtime-dynamics}
Within AIM, operational railway dynamics are interpreted as runtime
state evolution generated from railway runtime states.
\end{principle}

% ============================================================
\section{Dynamic Runtime States}
% ============================================================

\begin{definition}[Dynamic runtime state]
\label{def:runtime-dynamic-state}
A dynamic runtime state is an operational state generated from railway
runtime states during runtime interpretation.
\end{definition}

Dynamic runtime states do not belong to the intrinsic alignment model.
They are derived states.

They are generated from railway state, operational context, and runtime
interpretation.

Conceptually, a dynamic runtime state may be written as

\[
x(s)
=
\mathcal D
\bigl(
\mathrm{pose3}(s),
\mathcal C
\bigr),
\]

where \(\mathrm{pose3}(s)\) denotes the railway runtime state,
\(\mathcal C\) denotes operational context, and \(x(s)\) denotes the
resulting operational state.

% ============================================================
\section{Runtime Dynamics Operator}
% ============================================================

\begin{definition}[Runtime dynamics operator]
\label{def:runtime-dynamics-operator}
The runtime dynamics operator is the operator

\[
\mathcal D
\]

that maps railway runtime states and operational context to dynamic
runtime states.
\end{definition}

The operator \(\mathcal D\) represents the dynamic interpretation layer
of AIM.

It does not prescribe a particular vehicle model.
Instead, it describes where such models, simplified operational
relations, admissibility checks, and runtime evaluations are coupled to
the railway state.

Thus,

\[
\mathcal D
:
(\mathrm{pose3}, \mathcal C)
\mapsto
x
\]

expresses the transition from railway state to operational state.

% ============================================================
\section{Arc-Length State Evolution}
% ============================================================

\begin{definition}[Runtime state evolution]
\label{def:runtime-dynamics-state-evolution}
Runtime state evolution describes the change of a dynamic runtime state
along the alignment parameter.
\end{definition}

Within AIM, the natural evolution parameter is arc-length.

An abstract runtime evolution law may be written as

\[
x'(s)
=
f(x(s),u(s)),
\]

where \(x(s)\) denotes the dynamic runtime state, \(u(s)\) denotes
operational input, and \(f\) denotes the evolution law.

If a velocity function \(v(s)\) is known, time-based and
arc-length-based derivatives are connected by

\[
\frac{\dd}{\dd t}
=
v(s)
\frac{\dd}{\dd s}.
\]

This relation links intrinsic railway geometry to operational runtime
interpretation.

% ============================================================
\section{Geometry and Runtime Behaviour}
% ============================================================

Intrinsic geometry describes the structure of the railway.

Runtime dynamics describe the operational behaviour generated from that
structure.

The distinction is essential.

Curvature, profile, cant, and frame reconstruction define railway state.
Dynamic runtime states are derived from this state through operational
interpretation.

\begin{proposition}
\label{prop:dynamics-not-geometry}
Dynamic runtime states are derived from railway geometry, but they are
not themselves intrinsic geometric quantities.
\end{proposition}

% ============================================================
\section{Curvature as Runtime Generator}
% ============================================================

Runtime dynamics originate from intrinsic railway structure.

In particular, curvature evolution acts as one of the primary generators
of operational behaviour:

\[
\kappa(s)
\quad
\longrightarrow
\quad
\mathcal D
\quad
\longrightarrow
\quad
x(s).
\]

The quantities

\[
\kappa(s),
\qquad
\kappa'(s),
\qquad
\kappa''(s)
\]

are therefore not only geometric descriptors.
They influence the dynamic runtime states generated from the railway.

This observation is one of the central motivations of AIM:
operational behaviour is rooted in intrinsic curvature structure.

% ============================================================
\section{Runtime Dynamics and Vehicle Models}
% ============================================================

Runtime dynamics in AIM are not identical to full vehicle mechanics.

A complete vehicle-dynamics simulation may be attached to the AIM
runtime layer, but it is not required for defining the layer itself.

The role of runtime dynamics is to provide the operational
state-evolution interface between railway state and dynamic
interpretation.

\begin{principle}[Mechanics separation principle]
\label{princ:mechanics-separation}
Runtime dynamics in AIM form an operational state-evolution layer rather
than a complete vehicle-dynamics simulation layer.
\end{principle}

% ============================================================
\section{Canonical Interpretation}
% ============================================================

\begin{proposition}
Within AIM, railway dynamics are interpreted as runtime state evolution.
\end{proposition}

\begin{proposition}
Dynamic runtime states are generated from railway runtime states.
\end{proposition}

\begin{proposition}
The runtime dynamics operator \(\mathcal D\) maps railway runtime states
to operational state trajectories.
\end{proposition}

\begin{proposition}
Runtime dynamics belong to the operational interpretation layer and not
to intrinsic geometry itself.
\end{proposition}

\begin{proposition}
The AIM runtime-dynamics framework provides operational state evolution
rather than complete vehicle mechanics.
\end{proposition}

% ============================================================
\section{Connection to AIM}
% ============================================================

\begin{summary}
Runtime dynamics extend AIM from railway-state reconstruction toward
operational state evolution.

Within this framework:
\begin{itemize}
\item \(\mathrm{pose3}\) provides railway runtime states,
\item \(\mathcal D\) generates dynamic runtime states,
\item state evolution is naturally formulated over arc-length,
\item curvature evolution acts as a generator of runtime behaviour,
\item geometry and dynamics remain conceptually separated,
\item and detailed vehicle mechanics may be coupled through dedicated
runtime models.
\end{itemize}

Runtime dynamics therefore establish the bridge between intrinsic
geometry, railway state, vehicle-space interpretation, realization-aware
dynamics, and optimization.
\end{summary}
